\begin{pastexam}[label={px:ncu-111-dsa-12}]{中央資工所 111 學年度 多選題第 12 題}
（多選；每題 5 分，每個選項答對得 1 分，答錯倒扣 1 分）

Choose the correct statement(s):
\begin{enumerate}[label=(\Alph*)]
\item It has been proved that there exists no deterministic algorithm that can solve any given NPC
problem with a polynomial time complexity in the worst case.
\item If we can prove that the worst-case problem lower bound of an NPC problem is of a polynomial
order, then we can prove that NP$\neq$P.
\item If we can prove that an existing NP-hard problem is polynomially reducible to a problem X,
then X is NP-hard.
\item If a problem is proved to be an NP-hard and NP problem, then it is an NPC problem.
\item If a problem X can be solved by a non-deterministic algorithm taking a polynomial time
complexity in the worst case, then X is an NP problem.
\end{enumerate}
\end{pastexam}

\begin{pastsoln}
答案是 (C)、(D)、(E)。

名詞照資結書第~28~章：NP 是非決定性演算法能在多項式時間解的問題（等價地，答案是「是」時有多項式時間
可驗證的證據）；NP-hard 是每個 NP 問題都能多項式時間歸約過去的問題；NPC（NP 完備）是既屬於 NP 又是
NP-hard 的問題。
\begin{itemize}
\item (A) 錯。 這等於說已經證明了 $\mathrm{P} \ne \mathrm{NP}$，但這仍是未解的問題（資結書第~28.4~節）。
\item (B) 錯。 證明某個 NPC 問題的下界是多項式（例如 $\Omega(n^2)$），並不排除它有多項式時間的演算法，
推不出 $\mathrm{NP} \ne \mathrm{P}$。 要證明 $\mathrm{NP} \ne \mathrm{P}$，需要的是\textbf{超多項式}的
下界。
\item (C) 對。 設 $Y$ 是 NP-hard，$Y \le_p X$。 對任何 NP 問題 $Z$，$Z \le_p Y$，再接上 $Y \le_p X$；
多項式時間歸約可以串接（資結書引理~28.3），所以 $Z \le_p X$，$X$ 是 NP-hard。 這正是證明 NP-hard 的
標準做法（資結書第~28.4~節）。
\item (D) 對。 這就是 NPC 的定義：$X \in \mathrm{NP}$ 且 $X$ 是 NP-hard。
\item (E) 對。 這就是 NP 的定義（NP 的 N 是 nondeterministic）：存在非決定性演算法在最壞情況多項式時間
內解決它。
\end{itemize}
\end{pastsoln}
